\documentclass[10pt,a4paper]{amsart}
\usepackage{amsmath,amssymb,amsthm,graphicx}
% For commutative diagrams we use Xy-pic
\usepackage[all]{xy}

% We use 2cell for 2-commutative diagrams.
\xyoption{2cell}
% Theorem environments.
%
\theoremstyle{plain}
\newtheorem{theorem}[subsection]{Theorem}
\newtheorem{proposition}[subsection]{Proposition}
\newtheorem{lemma}[subsection]{Lemma}

\theoremstyle{definition}
\newtheorem{definition}[subsection]{Definition}
\newtheorem{example}[subsection]{Example}
\newtheorem{exercise}[subsection]{Exercise}
\newtheorem{situation}[subsection]{Situation}

\theoremstyle{remark}
\newtheorem{remark}[subsection]{Remark}
\newtheorem{remarks}[subsection]{Remarks}

\numberwithin{equation}{subsection}

% Macros
%
\def\lim{\mathop{\mathrm{lim}}\nolimits}
\def\colim{\mathop{\mathrm{colim}}\nolimits}
\def\Spec{\mathop{\mathrm{Spec}}}
\def\Hom{\mathop{\mathrm{Hom}}\nolimits}
\def\Ext{\mathop{\mathrm{Ext}}\nolimits}
\def\SheafHom{\mathop{\mathcal{H}\!\mathit{om}}\nolimits}
\def\SheafExt{\mathop{\mathcal{E}\!\mathit{xt}}\nolimits}
\def\Sch{\mathit{Sch}}
\def\Mor{\mathop{\mathrm{Mor}}\nolimits}
\def\Ob{\mathop{\mathrm{Ob}}\nolimits}
\def\Sh{\mathop{\mathit{Sh}}\nolimits}
\def\NL{\mathop{N\!L}\nolimits}
\def\CH{\mathop{\mathrm{CH}}\nolimits}
\def\proetale{{pro\text{-}\acute{e}tale}}
\def\etale{{\acute{e}tale}}
\def\QCoh{\mathit{QCoh}}
\def\Ker{\mathop{\mathrm{Ker}}}
\def\Im{\mathop{\mathrm{Im}}}
\def\Coker{\mathop{\mathrm{Coker}}}
\def\Coim{\mathop{\mathrm{Coim}}}

% Boxtimes
%
\DeclareMathSymbol{\boxtimes}{\mathbin}{AMSa}{"02}

%
% Macros for moduli stacks/spaces
%
\def\QCohstack{\mathcal{QC}\!\mathit{oh}}
\def\Cohstack{\mathcal{C}\!\mathit{oh}}
\def\Spacesstack{\mathcal{S}\!\mathit{paces}}
\def\Quotfunctor{\mathrm{Quot}}
\def\Hilbfunctor{\mathrm{Hilb}}
\def\Curvesstack{\mathcal{C}\!\mathit{urves}}
\def\Polarizedstack{\mathcal{P}\!\mathit{olarized}}
\def\Complexesstack{\mathcal{C}\!\mathit{omplexes}}
% \Pic is the operator that assigns to X its picard group, usage \Pic(X)
% \Picardstack_{X/B} denotes the Picard stack of X over B
% \Picardfunctor_{X/B} denotes the Picard functor of X over B
\def\Pic{\mathop{\mathrm{Pic}}\nolimits}
\def\Picardstack{\mathcal{P}\!\mathit{ic}}
\def\Picardfunctor{\mathrm{Pic}}
\def\Deformationcategory{\mathcal{D}\!\mathit{ef}}
\usepackage{cleveref}
\crefname{originaldefinition}{original definition}{original definitions}
\Crefname{originaldefinition}{Original definition}{Original definitions}
\makeatletter
\newlabel{definition-cycles}{{02QR}{source}{}{sourcepointer-1}{}}
\newlabel{definition-cycles@cref}{{[originaldefinition][0][]02QR}{[1][0][]source}{}{}{}}
\makeatother
\crefname{originallemma}{original lemma}{original lemmas}
\Crefname{originallemma}{Original lemma}{Original lemmas}
\makeatletter
\newlabel{lemma-flat-pullback-bivariant}{{0B78}{source}{}{sourcepointer-2}{}}
\newlabel{lemma-flat-pullback-bivariant@cref}{{[originallemma][0][]0B78}{[1][0][]source}{}{}{}}
\makeatother
\crefname{originallemma}{original lemma}{original lemmas}
\Crefname{originallemma}{Original lemma}{Original lemmas}
\makeatletter
\newlabel{lemma-push-proper-bivariant}{{0EPK}{source}{}{sourcepointer-3}{}}
\newlabel{lemma-push-proper-bivariant@cref}{{[originallemma][0][]0EPK}{[1][0][]source}{}{}{}}
\makeatother
\crefname{originallemma}{original lemma}{original lemmas}
\Crefname{originallemma}{Original lemma}{Original lemmas}
\makeatletter
\newlabel{lemma-bivariant-zero}{{02UC}{source}{}{sourcepointer-4}{}}
\newlabel{lemma-bivariant-zero@cref}{{[originallemma][0][]02UC}{[1][0][]source}{}{}{}}
\makeatother
\title{Bivariant Chow cohomology of a scheme of finite type over a field with target the point is ordinary Chow homology in complementary degree}
\author{The Stacks Project Authors}
\date{stacks-project:chow.tex:local-snapshot-20261012; GNU FDL 1.2 or later}
\begin{document}
\maketitle
\noindent\textit{Editorial note.} All delivered mathematical text is the author's own, copied from the pinned Stacks Project chapter file chow.tex; no statement or proof step was rewritten, reordered or removed.\par\medskip
\noindent\textit{Editorial note.} Preamble provenance: the pinned chapter file opens with an input of the shared preamble and carries no preamble text of its own, so the project's single Stacks-wide preamble staged at sources/stacks-chow/preamble.tex is excerpted verbatim as two spans, lines 14-18 (Xy-pic setup) and lines 154-217 (the native theorem declarations and mathematical macros). Nothing was added, removed or redefined; the wrapper class and the amsmath/amssymb/amsthm/graphicx base packages are supplied by the assembler.\par\medskip
\noindent\textit{Editorial note.} The counted statement and proof contain no cross reference, no citation and no included file, so the delivered body is self contained.\par\medskip
\noindent\textit{Editorial note.} Canonical location of the counted proposition: Stacks tag 0FBW (sources/stacks-chow/pinned-original-source.tex line 14183).\par\medskip
\noindent\textit{Allowed standard prerequisites (editorial).}\begin{itemize}
\item Bivariant operational Chow groups A\textasciicircum{}p(X -> Y) and their axioms, including the empty-class reduction. Same chapter, chow.tex, Lemma labelled lemma-bivariant-zero (source line 6585), invoked at proof line 14236. Supplies that it suffices to test equality of bivariant classes against the classes of integral schemes; the construction of the class is inside the counted span.
\item Flat pullback and proper pushforward are bivariant operations. Same chapter, Lemmas labelled lemma-flat-pullback-bivariant (source line 6131) and lemma-push-proper-bivariant (source line 6167), invoked at proof lines 14214 and 14222. Supply the two operations whose composite g\_* o nu o f\textasciicircum{}* realises the desired class; the identification of that composite with c\_alpha and the inverse check are in the counted span.
\item Writing a cycle class as a finite sum of integral closed subschemes. Same chapter, Definition labelled definition-cycles (source line 1535), invoked at proof line 14206. Supplies the representation alpha = sum n\_i [X\_i] that the construction of the bivariant class starts from.
\end{itemize}\medskip
\ifdefined\hypertarget\hypertarget{sourcepointer-1}{}\fi
\noindent\textit{Source reference.} Original definition 02QR: sources/stacks-chow/pinned-original-source.tex, source lines 1535--1535. Same-chapter pointer into the author environment definition at source line 1535; invoked at proof line 14206 of the counted proof. It is retained only as a navigation pointer; every load-bearing step is inside the delivered span.\par\medskip
\ifdefined\hypertarget\hypertarget{sourcepointer-2}{}\fi
\noindent\textit{Source reference.} Original lemma 0B78: sources/stacks-chow/pinned-original-source.tex, source lines 6131--6131. Same-chapter pointer into the author environment lemma at source line 6131; invoked at proof line 14214 of the counted proof. It is retained only as a navigation pointer; every load-bearing step is inside the delivered span.\par\medskip
\ifdefined\hypertarget\hypertarget{sourcepointer-3}{}\fi
\noindent\textit{Source reference.} Original lemma 0EPK: sources/stacks-chow/pinned-original-source.tex, source lines 6167--6167. Same-chapter pointer into the author environment lemma at source line 6167; invoked at proof line 14222 of the counted proof. It is retained only as a navigation pointer; every load-bearing step is inside the delivered span.\par\medskip
\ifdefined\hypertarget\hypertarget{sourcepointer-4}{}\fi
\noindent\textit{Source reference.} Original lemma 02UC: sources/stacks-chow/pinned-original-source.tex, source lines 6585--6585. Same-chapter pointer into the author environment lemma at source line 6585; invoked at proof line 14236 of the counted proof. It is retained only as a navigation pointer; every load-bearing step is inside the delivered span.\par\medskip
\begin{lemma}
\label{lemma-chow-cohomology-towards-point}
Let $k$ be a field. Let $X$ be a scheme locally of finite type over $k$.
Then we have a canonical identification
$$
A^p(X \to \Spec(k)) = \CH_{-p}(X)
$$
for all $p \in \mathbf{Z}$.
\end{lemma}

\begin{proof}
Consider the element $[\Spec(k)] \in \CH_0(\Spec(k))$. We get a map
$A^p(X \to \Spec(k)) \to \CH_{-p}(X)$ by sending $c$ to $c \cap [\Spec(k)]$.

\medskip\noindent
Conversely, suppose we have $\alpha \in \CH_{-p}(X)$.
Then we can define $c_\alpha \in A^p(X \to \Spec(k))$ as
follows: given $X' \to \Spec(k)$ and $\alpha' \in \CH_n(X')$
we let
$$
c_\alpha \cap \alpha' = \alpha \times \alpha'
$$
in $\CH_{n - p}(X \times_k X')$. To show that this is a bivariant
class we write $\alpha = \sum n_i[X_i]$ as in
Definition \ref{definition-cycles}. Consider the composition
$$
\coprod X_i \xrightarrow{g} X \to \Spec(k)
$$
and denote $f : \coprod X_i \to \Spec(k)$ the composition.
Then $g$ is proper and $f$ is flat of relative dimension $-p$.
Pullback along $f$ is a bivariant class
$f^* \in A^p(\coprod X_i \to \Spec(k))$ by
Lemma \ref{lemma-flat-pullback-bivariant}.
Denote $\nu \in A^0(\coprod X_i)$ the bivariant class
which multiplies a cycle by $n_i$ on the $i$th component.
Thus $\nu \circ f^* \in A^p(\coprod X_i \to X)$.
Finally, we have a bivariant class
$$
g_* \circ \nu \circ f^*
$$
by Lemma \ref{lemma-push-proper-bivariant}. The reader easily
verifies that $c_\alpha$ is equal to this class and hence
is itself a bivariant class.

\medskip\noindent
To finish the proof we have to show that the two constructions
are mutually inverse. Since $c_\alpha \cap [\Spec(k)] = \alpha$
this is clear for one of the two directions. For the other, let
$c \in A^p(X \to \Spec(k))$ and set $\alpha = c \cap [\Spec(k)]$.
It suffices to prove that
$$
c \cap [X'] = c_\alpha \cap [X']
$$
when $X'$ is an integral scheme locally of finite type over $\Spec(k)$,
see Lemma \ref{lemma-bivariant-zero}. However, then $p' : X' \to \Spec(k)$
is flat of relative dimension $\dim(X')$ and hence
$[X'] = (p')^*[\Spec(k)]$. Thus the fact that the bivariant classes
$c$ and $c_\alpha$ agree on $[\Spec(k)]$ implies they
agree when capped against $[X']$ and the proof is complete.
\end{proof}

\end{document}
